% HOST-INDEPENDENT PREAMBLE (added by scripts/fix_insert_portability.py).
% This insert may be \input by either root manuscript, and they do not define the
% same environments or macros: w33_paper.tex has lemma/proposition/remark and
% \PSp/\Aut, photonic_holonet.tex has none of them. Define only what is missing,
% so the file is portable and re-running this is a no-op.
\makeatletter
\@ifundefined{theorem}{\newtheorem{theorem}{Theorem}}{}
\@ifundefined{lemma}{\newtheorem{lemma}[theorem]{Lemma}}{}
\@ifundefined{proposition}{\newtheorem{proposition}[theorem]{Proposition}}{}
\@ifundefined{remark}{\newtheorem{remark}[theorem]{Remark}}{}
\@ifundefined{corollary}{\newtheorem{corollary}[theorem]{Corollary}}{}
\@ifundefined{definition}{\newtheorem{definition}[theorem]{Definition}}{}
\makeatother
\providecommand{\PSp}{\mathrm{PSp}}
\providecommand{\Aut}{\mathrm{Aut}}
\providecommand{\spec}{\operatorname{spec}}


\section{The selector-matching association scheme}
\label{sec:bt1355-bt1359-selector-matching-scheme}

Pass~1353 identifies the $120$ literal length-four selectors with pairs
$(L,M)$, where $L$ is a totally isotropic line of $W(3,3)$ and $M$ is
one of the three perfect matchings of the four points of $L$.  The
following result determines the complete transport algebra on those
selectors.

\begin{theorem}[Selector-matching association scheme]
Let
\[
 X=\{(L,M):L\text{ is a line of }W(3,3),\ M\text{ is a perfect
 matching of }L\}.
\]
For two members $(L,M),(L',M')\in X$, let $R_0$ be equality, $R_1$ mean
the same line with different matchings, and $R_2$ mean distinct
intersecting lines.  If $L$ and $L'$ are disjoint, the generalized-
quadrangle axiom gives a unique transversal bijection $\tau_{L,L'}:L\to
L'$.  Let $R_3$ mean that $\tau_{L,L'}(M)=M'$, and let $R_4$ be the
remaining disjoint-line case.

Then $\{R_0,R_1,R_2,R_3,R_4\}$ is a symmetric commutative four-class
association scheme on $120$ points.  Its valencies and primitive
multiplicities are, respectively,
\[
 \boxed{(1,2,36,27,54)},
 \qquad
 \boxed{(1,15,24,20,60)}.
\]
With the relation ordering displayed above, its first and second
eigenmatrices are
\[
P=\begin{pmatrix}
1&2&36&27&54\\
1&2&-12&3&6\\
1&2&6&-3&-6\\
1&-1&0&9&-9\\
1&-1&0&-3&3
\end{pmatrix},
\]
\[
Q=\begin{pmatrix}
1&15&24&20&60\\
1&15&24&-10&-30\\
1&-5&4&0&0\\
1&5/3&-8/3&20/3&-20/3\\
1&5/3&-8/3&-10/3&10/3
\end{pmatrix},
\qquad PQ=QP=120I_5.
\]
The parabolic relation $R_0\cup R_1$ has forty fibers of size three,
and the quotient is the line-intersection graph
$\operatorname{SRG}(40,12,2,4)$.  The fiber-constant primitive ranks are
$1,15,24$, while the fiber-zero ranks are $20,60$.
\end{theorem}

\begin{proof}
Construct $W(3,3)$ as the projective isotropic points of the standard
symplectic form on $\mathbb F_3^4$.  Its forty maximal isotropic
$4$-subsets are the lines.  For disjoint lines, each point of the first
is collinear with exactly one point of the second, giving the stated
transversal bijection.  Direct exact enumeration shows that the five
relations partition $X^2$, are symmetric, and have constant
intersection numbers.  The displayed rows of $P$ satisfy every
character equation of this intersection algebra; inversion gives $Q$
and $PQ=QP=120I_5$.  The same certificate checks every fusion and every
P- and Q-polynomial ordering.
\end{proof}

\begin{corollary}[Fusion rigidity and transport holonomy]
Among the fifteen possible partitions of the four nonidentity
relations, the only proper fusions are obtained by fusing $R_3$ with
$R_4$, then optionally fusing all three inter-line relations, or fusing
all nonidentity relations.  The full scheme is neither P-polynomial nor
Q-polynomial.

The disjoint-line graph is connected, $27$-regular, and has $540$
edges.  Its transversal transports act on each three-element matching
fiber and generate the full group $S_3$.  Since the centralizer of this
holonomy in $S_3$ is trivial, the kernel of the action of
$\operatorname{Aut}(\mathcal X)$ on the forty fibers is trivial.  Hence
\[
 \operatorname{Aut}(\mathcal X)\cong\operatorname{PGSp}(4,3)
 \cong W(E_6),
 \qquad |\operatorname{Aut}(\mathcal X)|=51840.
\]
\end{corollary}

\paragraph{Scope boundary.}
This closes the previously named selector-transport algebra: the
$2,27,36,54$ orbitals are the relations of one exact imprimitive
association scheme, and their line-dependent connection has full
$S_3$ holonomy.  It does not produce a $12$-regular $H_4$/600-cell
adjacency, select a preferred matching in any fiber, or supply evidence
for a physical interpretation beyond the finite transport model.
